\section{Perturbative stability of nearly spherical sets}
\label{sec:nearly}

This section proves, in full and self-containedly, the perturbative
torsional stability of small-amplitude radial graphs consumed by the
reduction of \cref{sec:main}. The result is a Fuglede-type estimate
\emph{for the torsion functional}: a planar set whose boundary is a graph of
small amplitude over a circle, after the natural volume and barycentre
normalisations, has scale-invariant torsional deficit controlling the square
of its Fraenkel asymmetry. We invoke no selection principle, no
spectral black box, and no quantitative isoperimetric inequality. The proof
is an explicit second-order torsion computation through the dual
(complementary) variational principle, with an admissible vector field built
by hand; its only nonclassical ingredient is a frequency-dependent damping of
the harmonic corrector that controls high-frequency boundary oscillations and
lets the hypothesis be merely $\|\varphi\|_{L^\infty}\le\eta_G$.

\subsection{Statement and standing normalisation}

Throughout, $S^1=\R/2\pi\Z$, $\omega(\theta)=(\cos\theta,\sin\theta)$, and for
$\varphi\in \mathrm{Lip}(S^1)$ with $\|\varphi\|_{L^\infty}<1$ and $z\in\R^2$
we write
\begin{equation}\label{eq:Gdef}
   G=G(z,\hat r_G,\varphi):=\set{z+r\omega(\theta):\ \theta\in S^1,\
   0\le r<\hat r_G\,(1+\varphi(\theta))}
\end{equation}
for the radial graph of amplitude profile $\varphi$, base radius
$\hat r_G=\sqrt{|G|/\pi}$ and centre $z$. We recall that, for a bounded set
$E$ of finite measure with torsion $T(E)=\int_E u_E$ ($-\Delta u_E=1$,
$u_E\in H^1_0(E)$),
\begin{equation}\label{eq:dstarAdef}
   \dstar(E)=\frac{T(B_E)-T(E)}{T(B_E)},\quad T(B_E)=\frac{|E|^2}{8\pi},
   \qquad
   \AF(E)=\inf_{x_0}\frac{|E\triangle B(x_0,r_E)|}{|E|},\ r_E=\sqrt{|E|/\pi},
\end{equation}
are the scale-invariant deficit and the Fraenkel asymmetry; both are
invariant under translations and dilations.

\begin{theorem}[Radial-graph torsion stability]\label{thm:graphG}
There are absolute constants
\begin{equation}\label{eq:graphG-const}
   \eta_G=\tfrac1{150},\qquad C_G=48,
\end{equation}
such that the following holds. Let $G=G(z,\hat r_G,\varphi)$ be a radial graph
as in \eqref{eq:Gdef} with $\varphi\in\mathrm{Lip}(S^1)$,
$\|\varphi\|_{L^\infty}\le\eta_G$, and barycentre at $z$. Then
\begin{equation}\label{eq:graphG}
   \AF(G)^2\ \le\ C_G\,\dstar(G).
\end{equation}
The statement is scale-invariant; $G$ need not have area $\pi$. It uses no
property of the original set $\Om$ and no selection.
\end{theorem}

Since $\dstar$ and $\AF$ are dilation- and translation-invariant, replacing
$G$ by $\hat r_G^{-1}(G-z)$ we assume for the proof
\begin{equation}\label{eq:norm}
   z=0,\qquad \hat r_G=1,\qquad\text{so}\qquad |G|=\pi,
\end{equation}
which fixes the volume constraint
\begin{equation}\label{eq:vol-constraint}
   |G|=\int_0^{2\pi}\frac{(1+\varphi)^2}{2}\d\theta=\pi
   \quad\Longleftrightarrow\quad
   \int_0^{2\pi}\bigl(2\varphi+\varphi^2\bigr)\d\theta=0 .
\end{equation}
Under \eqref{eq:norm}, $T(B_G)=\pi/8$, while $u_{B_1}=(1-|x|^2)/4$ and
$T(B_1)=\pi/8$. Expand $\varphi$ in its real Fourier series,
\begin{equation}\label{eq:fourier}
   \varphi(\theta)=a_0+\sum_{k\ge1}\bigl(a_k\cos k\theta+b_k\sin k\theta\bigr),
   \qquad
   \|\varphi\|_{L^2(S^1)}^2=\pi\Bigl(2a_0^2+\sum_{k\ge1}(a_k^2+b_k^2)\Bigr) .
\end{equation}
Lipschitz regularity gives $\sum_k k^2(a_k^2+b_k^2)<\infty$, so the series
below converge in the indicated norms.

\subsection{The dual (complementary) variational principle}

The torsion has two dual characterisations: the Dirichlet principle
$T(E)=\max_{w\in H^1_0(E)}\bigl(2\int_E w-\int_E|\nabla w|^2\bigr)$ (attained
at $w=u_E$), and the complementary principle below, which produces
\emph{upper} bounds on $T(E)$ — equivalently \emph{lower} bounds on the
deficit — from explicit divergence-constrained fields.

\begin{lemma}[Complementary principle]\label{lem:dual}
Let $E\subset\R^2$ be bounded open with finite perimeter and torsion function
$u_E\in H^1_0(E)$. For every $\sigma\in L^2(E;\R^2)$ with $\divg\sigma=-1$ in
$\mathcal D'(E)$,
\begin{equation}\label{eq:dual}
   T(E)\ \le\ \int_E|\sigma|^2,
\end{equation}
with equality iff $\sigma=\nabla u_E$.
\end{lemma}

\begin{proof}
Since $u_E\in H^1_0(E)$ and $\divg\sigma=-1$, the pairing
$\int_E\sigma\cdot\nabla u_E=-\langle\divg\sigma,u_E\rangle=\int_E u_E$ is
well defined, and $\int_E|\nabla u_E|^2=\int_E u_E=T(E)$ by testing
$-\Delta u_E=1$ with $u_E$. Hence
\[
   0\le\int_E|\sigma-\nabla u_E|^2
   =\int_E|\sigma|^2-2\!\int_E\sigma\cdot\nabla u_E+\!\int_E|\nabla u_E|^2
   =\int_E|\sigma|^2-T(E),
\]
which is \eqref{eq:dual}; equality forces $\sigma=\nabla u_E$ a.e.
\end{proof}

The disc field $\sigma_0:=\nabla u_{B_1}=-x/2$ satisfies $\divg\sigma_0=-1$
and $\int_{B_1}|\sigma_0|^2=T(B_1)=\pi/8$. We perturb it into an admissible
field on $G$.

\subsection{The comparison field}

Set the corrector coefficients
\begin{equation}\label{eq:gk}
   \alpha_k:=\tfrac12 m_k\,a_k,\quad \beta_k:=\tfrac12 m_k\,b_k,
   \qquad m_k:=\min\Bigl(1,\tfrac2k\Bigr)\in(0,1],\quad k\ge1 .
\end{equation}
The damping $m_k$ equals $1$ at the critical mode $k=2$ and decays like $2/k$
for $k>2$, so the corrector \emph{slope} stays bounded:
\begin{equation}\label{eq:slopebound}
   k|\alpha_k|=\tfrac12 m_k k\,|a_k|\le|a_k|,
   \qquad k|\beta_k|\le|b_k| .
\end{equation}
Define the harmonic corrector and its normal trace on $\partial B_1$,
\begin{equation}\label{eq:hC}
   h(r,\theta):=\sum_{k\ge1}r^k\bigl(\alpha_k\cos k\theta+\beta_k\sin k\theta\bigr),
   \qquad
   c(\theta):=\partial_r h(1,\theta)=\sum_{k\ge1}k\bigl(\alpha_k\cos k\theta+\beta_k\sin k\theta\bigr).
\end{equation}
By \eqref{eq:slopebound} and Lipschitz regularity, $\nabla h$ and $c$ are in
$L^2$, $\Delta h=0$ in $\R^2$, and
\begin{equation}\label{eq:cL2}
   \|c\|_{L^2(S^1)}^2=\pi\sum_{k\ge1}k^2(\alpha_k^2+\beta_k^2)
   \le\pi\sum_{k\ge1}(a_k^2+b_k^2)\le\|\varphi\|_{L^2}^2 .
\end{equation}
The comparison field on $G$ is
\begin{equation}\label{eq:sigma}
   \sigma(x):=
   \begin{cases}
      \nabla\bigl(u_{B_1}+h\bigr)(x), & x\in G,\ |x|\le1,\\[1mm]
      \Bigl(-\dfrac r2+\dfrac{c(\theta)}{r}\Bigr)\,\omega(\theta),
        & x=r\omega(\theta)\in G,\ |x|>1 .
   \end{cases}
\end{equation}
The collar line is used only where $\varphi(\theta)>0$, i.e.\ where $G$
protrudes beyond $B_1$.

\begin{lemma}[Admissibility]\label{lem:admissible}
The field $\sigma$ of \eqref{eq:sigma} lies in $L^2(G;\R^2)$ and satisfies
$\divg\sigma=-1$ in $\mathcal D'(G)$.
\end{lemma}

\begin{proof}
On $\{|x|<1\}$, $\divg\sigma=\Delta(u_{B_1}+h)=-1+0=-1$. On the collar
$\{|x|>1\}$, for a radial field $f(r)\omega$ one has
$\divg(f\omega)=\frac1r\frac{d}{dr}(rf)$; with $f=-\frac r2+\frac{c}{r}$ this
is $\frac1r\frac{d}{dr}\bigl(-\frac{r^2}2+c\bigr)=-1$. Across the interface
$\Gamma=\partial B_1\cap G$, the normal (radial) components match: from inside
$\sigma\cdot\omega=\partial_r(u_{B_1}+h)|_{r=1}=-\tfrac12+c(\theta)$, from the
collar $\sigma\cdot\omega=-\tfrac12+c(\theta)$. Hence the distributional
divergence carries no surface measure on $\Gamma$ and equals $-1$ on all of
$G$. Square-integrability: $\nabla u_{B_1}$ is smooth, $\nabla h\in
L^2(B_{1+\eta_G})$ by \eqref{eq:slopebound}, and the collar is a thin set of
width $\le\eta_G$ on which $|\sigma|^2\le2(\tfrac r2)^2+2c(\theta)^2$.
\end{proof}

\subsection{The exact second-order energy}

We compute $\int_G|\sigma|^2$ exactly and separate the quadratic form from a
genuinely cubic remainder. Write $p:=\varphi(\theta)$ and
\begin{equation}\label{eq:Phidef}
   \Phi_{\mathrm{col}}(r,\theta):=\Bigl(-\tfrac r2+\tfrac{c}{r}\Bigr)^2 r,
   \qquad
   \Phi_{\mathrm{in}}(r,\theta):=\Bigl[\bigl(-\tfrac r2+\partial_r h\bigr)^2
   +\bigl(\tfrac1r\partial_\theta h\bigr)^2\Bigr]r ,
\end{equation}
the polar energy densities of the collar and interior fields.

\begin{lemma}[Energy identity]\label{lem:energy}
With $\sigma$ as in \eqref{eq:sigma},
\begin{equation}\label{eq:energy-split}
   \int_G|\sigma|^2
   =\frac\pi8
     +\pi\sum_{k\ge1}k(\alpha_k^2+\beta_k^2)
     +\mathcal N,
   \qquad
   \mathcal N=\int_{\{p>0\}}\!\!\int_1^{1+p}\!\!\Phi_{\mathrm{col}}\,dr\,d\theta
     -\int_{\{p<0\}}\!\!\int_{1+p}^1\!\!\Phi_{\mathrm{in}}\,dr\,d\theta .
\end{equation}
\end{lemma}

\begin{proof}
Split $\int_G=\int_{B_1}-\int_{B_1\setminus G}+\int_{G\setminus B_1}$. Over
the full disc, $\int_{B_1}|\nabla(u_{B_1}+h)|^2
=\int_{B_1}|\nabla u_{B_1}|^2+2\int_{B_1}\nabla u_{B_1}\cdot\nabla
h+\int_{B_1}|\nabla h|^2$. The first term is $\pi/8$; the cross term vanishes,
$\int_{B_1}\nabla u_{B_1}\cdot\nabla h=-\int_{B_1}u_{B_1}\Delta
h+\int_{\partial B_1}u_{B_1}\partial_\nu h=0$ (as $\Delta h=0$,
$u_{B_1}|_{\partial B_1}=0$); and, by orthogonality of the harmonic modes,
$\int_{B_1}|\nabla h|^2=\pi\sum_k k(\alpha_k^2+\beta_k^2)$ since
$\int_{B_1}|\nabla(r^k\cos k\theta)|^2=\pi k$. The dimple
$-\int_{B_1\setminus G}|\nabla(u_{B_1}+h)|^2$ (where $p<0$) and the collar
$+\int_{G\setminus B_1}|\sigma|^2$ (where $p>0$) are precisely $\mathcal N$.
\end{proof}

\begin{lemma}[Quadratic form and cubic remainder]\label{lem:BL}
Assume $\|\varphi\|_{L^\infty}\le\tfrac1{150}$ and the corrector
\eqref{eq:gk}. Then
\begin{equation}\label{eq:quadform}
   \int_G|\sigma|^2
   =\frac\pi8
     +\frac14\!\int_{S^1}\!\varphi
     +\frac38\!\int_{S^1}\!\varphi^2
     +\pi\!\sum_{k\ge1}\!\bigl(k\alpha_k^2-k\alpha_k a_k\bigr)
     +\pi\!\sum_{k\ge1}\!\bigl(k\beta_k^2-k\beta_k b_k\bigr)
     +\mathcal R,
\end{equation}
where the remainder obeys the explicit bound
\begin{equation}\label{eq:Rbound}
   |\mathcal R|\ \le\ 5\,\|\varphi\|_{L^\infty}\bigl(\|\varphi\|_{L^2}^2+\|c\|_{L^2}^2\bigr)
   \ \le\ 10\,\|\varphi\|_{L^\infty}\,\|\varphi\|_{L^2}^2 ,
\end{equation}
the last step by \eqref{eq:cL2}.
\end{lemma}

\begin{proof}
We expand $\mathcal N$ from \eqref{eq:energy-split} to second order in $p$ and
collect the cubic tail into $\mathcal R$. The two boundary densities are
$C^\infty$ in $r$ on $[\tfrac12,2]\supset[1-\eta_G,1+\eta_G]$, with values and
$r$-derivatives at $r=1$
\begin{equation}\label{eq:traces}
   \Phi_{\mathrm{col}}(1)=\bigl(\tfrac12-c\bigr)^2,\quad
   \Phi_{\mathrm{in}}(1)=\bigl(\tfrac12-c\bigr)^2+s^2,\quad
   \Phi_{\mathrm{col}}'(1)=\tfrac34-c-c^2,
\end{equation}
where $s(\theta):=\tfrac1r\partial_\theta h(1,\theta)$,
$\|s\|_{L^2}=\|c\|_{L^2}$ (the tangential and normal traces of a harmonic
function on $\partial B_1$ have equal $L^2$ norm). For each $\theta$, Taylor's
theorem with integral remainder gives, on the collar branch $p>0$,
\[
   \int_1^{1+p}\!\Phi_{\mathrm{col}}\,dr
   =p\,\Phi_{\mathrm{col}}(1)+\tfrac12 p^2\,\Phi_{\mathrm{col}}'(1)
   +\int_1^{1+p}\tfrac12(1+p-r)^2\,\Phi_{\mathrm{col}}''(r)\,dr,
\]
and analogously on the dimple branch $p<0$ for $\Phi_{\mathrm{in}}$. We keep
the \emph{disc} part of the first two terms in the quadratic form and put the
$c,s$-dependent part and the third-order tail into $\mathcal R$. Concretely,
writing $\Phi_{\mathrm{col}}(1)=\tfrac14-c+c^2$ and
$\tfrac12\Phi_{\mathrm{col}}'(1)=\tfrac38-\tfrac c2-\tfrac{c^2}2$, the kept
quadratic contribution of the collar branch is
$\int_{\{p>0\}}\bigl(\tfrac14 p-pc+\tfrac38 p^2\bigr)$, and the residue
$\int_{\{p>0\}}\bigl(p\,c^2-\tfrac12 p^2(c+c^2)\bigr)+(\text{tail})$ goes to
$\mathcal R$; the dimple branch is identical except for the extra
$\int_{\{p<0\}}p\,s^2$ (from $\Phi_{\mathrm{in}}(1)$), also placed in
$\mathcal R$. Since $\Phi_{\mathrm{col}}(1)$ and $\Phi_{\mathrm{in}}(1)$ share
the kept part $\tfrac14-c$, the kept quadratic over the full circle is
\[
   \int_{S^1}\Bigl(\tfrac14\varphi-\varphi\,c+\tfrac38\varphi^2\Bigr)
   =\tfrac14\!\int\!\varphi+\tfrac38\!\int\!\varphi^2
     -\pi\!\sum_k k(\alpha_k a_k+\beta_k b_k),
\]
the last equality by \eqref{eq:hC} and orthogonality. Adding the harmonic
term $\pi\sum_k k(\alpha_k^2+\beta_k^2)$ from \eqref{eq:energy-split} gives
exactly the displayed quadratic form in \eqref{eq:quadform}.

\emph{Remainder bound.} Each residue is controlled pointwise by $|p|\le\eta_G$
and the traces $c,s$. The second-order $c$-residues satisfy
\[
   \Bigl|\int_{S^1}\!\bigl(\varphi\,c^2-\tfrac12\varphi^2(c+c^2)\bigr)\Bigr|
   +\Bigl|\int_{\{p<0\}}\!\varphi\,s^2\Bigr|
   \le \eta_G\!\int\!\bigl(c^2+s^2\bigr)+\tfrac12\eta_G^2\!\int\!\bigl(|c|+c^2\bigr)
   \le 2\eta_G\,\|c\|_{L^2}^2,
\]
using $\|s\|_{L^2}=\|c\|_{L^2}$, $\eta_G\le\tfrac1{150}$, and
$\int|c|\le\sqrt{2\pi}\|c\|_{L^2}\le\|c\|_{L^2}^2/\eta_G$ when
$\|c\|_{L^2}\ge\sqrt{2\pi}\,\eta_G$ (otherwise the term is itself
$\le2\eta_G\|c\|_{L^2}^2$ trivially). For the Taylor tails, on the relevant
$r$-interval $[1-\eta_G,1+\eta_G]$ one has the bounds
$|\Phi_{\mathrm{col}}''|\le 2(1+c^2)$ and $|\Phi_{\mathrm{in}}''|\le
2(1+c^2+s^2)$ (direct differentiation of \eqref{eq:Phidef}, using $r\le
1+\eta_G\le\tfrac{151}{150}$ and, on the dimple, $r\le1$ so $r^{k-1}\le1$,
whence the corrector contributes only through $c,s$ up to an absolute
factor), giving
\[
   |\text{tail}|\le\tfrac16\!\int|p|^3\bigl(1+c^2+s^2\bigr)
   \le\tfrac16\eta_G\!\int\varphi^2+\tfrac16\eta_G\!\int\varphi^2(c^2+s^2)
   \le\tfrac16\eta_G\,\|\varphi\|_{L^2}^2+\tfrac13\eta_G\,\|c\|_{L^2}^2 .
\]
Collecting, $|\mathcal R|\le\eta_G\bigl(\tfrac16\|\varphi\|_{L^2}^2+(2+\tfrac13)\|c\|_{L^2}^2\bigr)
\le 5\,\eta_G(\|\varphi\|_{L^2}^2+\|c\|_{L^2}^2)$, which is \eqref{eq:Rbound}.
\end{proof}

\begin{remark}[On the damping $m_k$]\label{rem:damping}
The undamped corrector $\alpha_k=\tfrac12 a_k$ minimises the quadratic form
mode by mode and gives the sharp coefficient $\tfrac\pi8(3-2k)$, but its
normal trace has $\|c\|_{L^2}^2=\tfrac\pi4\sum_k k^2(a_k^2+b_k^2)$, an
$H^1$-quantity \emph{not} controlled by $\|\varphi\|_{L^2}$; then the
remainder would carry an unbounded factor of $k$. The damping
\eqref{eq:gk} caps the slope, $\|c\|_{L^2}\le\|\varphi\|_{L^2}$ by
\eqref{eq:cL2}, so $\mathcal R$ is genuinely cubic \emph{uniformly in the
frequency $k$}, at the cost of only a bounded factor in the leading
coefficient (preserved exactly at the critical mode $k=2$). This is the sole
place high-frequency oscillations are handled, and it is what allows the
hypothesis to be merely $\|\varphi\|_{L^\infty}\le\eta_G$.
\end{remark}

\subsection{Coercivity of the deficit}

We first read off the per-mode quadratic coefficient.

\begin{lemma}[Per-mode deficit coefficient]\label{lem:permode}
Using the volume identity
$\tfrac14\int_{S^1}\varphi+\tfrac38\int_{S^1}\varphi^2=-\tfrac14\int_{S^1}\varphi^2$
(immediate from \eqref{eq:vol-constraint}, since
$\int\varphi=-\int\varphi^2$), the net coefficient of
$a_k^2$ in $\tfrac\pi8-\int_G|\sigma|^2$ (and likewise of $b_k^2$) is
\begin{equation}\label{eq:Dk}
   D_k:=\pi\Bigl(\tfrac12 k m_k-\tfrac14 k m_k^2-\tfrac14\Bigr),
   \qquad m_k=\min(1,2/k).
\end{equation}
For $k\ge2$ one has
\begin{equation}\label{eq:Dk-floor}
   D_2=\frac\pi4,\qquad D_k=\pi\Bigl(\tfrac34-\tfrac1k\Bigr)\ \ge\ \frac{5\pi}{12}
   \quad(k\ge3),
   \qquad\text{hence } D_k\ge\frac\pi4\ \ \forall k\ge2,
\end{equation}
while $D_1=0$.
\end{lemma}

\begin{proof}
By \eqref{eq:gk}, $k\alpha_k^2=\tfrac14 k m_k^2 a_k^2$ and $k\alpha_k
a_k=\tfrac12 k m_k a_k^2$, so the contribution of the $k$-th cosine mode to
the corrector sum in \eqref{eq:quadform} is $\pi(k\alpha_k^2-k\alpha_k
a_k)=\pi(\tfrac14 k m_k^2-\tfrac12 k m_k)a_k^2$. The disc-collar quadratic
$\tfrac38\int\varphi^2$ and the linear $\tfrac14\int\varphi$ combine, by the
volume identity above, into $-\tfrac14\int\varphi^2$, whose
$k$-th cosine contribution is $-\tfrac14\pi a_k^2$. Hence the coefficient of
$a_k^2$ in $\int_G|\sigma|^2-\tfrac\pi8$ is
$\pi(\tfrac14 k m_k^2-\tfrac12 k m_k+\tfrac14)$, and its negative is $D_k$ as
in \eqref{eq:Dk}. For $k=2$ ($m_2=1$): $D_2=\pi(1-\tfrac12-\tfrac14)=\tfrac\pi4$.
For $k\ge3$ ($m_k=2/k$): $\tfrac12 k m_k=1$, $\tfrac14 k m_k^2=\tfrac1k$, so
$D_k=\pi(1-\tfrac1k-\tfrac14)=\pi(\tfrac34-\tfrac1k)\ge\pi(\tfrac34-\tfrac13)=\tfrac{5\pi}{12}$.
For $k=1$ ($m_1=1$): $D_1=\pi(\tfrac12-\tfrac14-\tfrac14)=0$.
\end{proof}

\begin{lemma}[Normalisation kills modes $0$ and $1$]\label{lem:normmodes}
Under \eqref{eq:vol-constraint} and barycentre $0$, with
$\|\varphi\|_{L^\infty}\le\eta_G\le\tfrac1{150}$ and
$\rho^2:=\sum_{k\ge2}(a_k^2+b_k^2)$,
\begin{equation}\label{eq:lowmodes}
   2a_0^2+a_1^2+b_1^2\ \le\ \tfrac1{10}\,\rho^2,
   \qquad\text{hence}\qquad
   \|\varphi\|_{L^2}^2=\pi\bigl(2a_0^2+a_1^2+b_1^2+\rho^2\bigr)\le\tfrac{11}{10}\pi\,\rho^2 .
\end{equation}
\end{lemma}

\begin{proof}
By \eqref{eq:vol-constraint}, $2\pi a_0=-\int\varphi^2$, so
$|a_0|\le\tfrac1{2\pi}\|\varphi\|_{L^\infty}\!\int|\varphi|\le\tfrac1{2\pi}\eta_G\sqrt{2\pi}\,\|\varphi\|_{L^2}$,
whence $2a_0^2\le\tfrac{\eta_G^2}\pi\|\varphi\|_{L^2}^2$. For the barycentre,
$\int_G x\,dx=\int_0^{2\pi}\tfrac{(1+\varphi)^3}3\,\omega\,d\theta$; its linear
part is $\int\varphi\,\omega=\pi(a_1,b_1)$, and setting the moment to $0$ gives
$\pi^2(a_1^2+b_1^2)\le\bigl(\tfrac13\!\int(3|\varphi|^2+|\varphi|^3)\bigr)^2
\le(2\!\int\varphi^2)^2\le4\eta_G^2(\sqrt{2\pi}\|\varphi\|_{L^2})^2
=8\pi\eta_G^2\|\varphi\|_{L^2}^2$. Thus
$2a_0^2+a_1^2+b_1^2\le\bigl(\tfrac{\eta_G^2}\pi+\tfrac{8\eta_G^2}\pi\bigr)\|\varphi\|_{L^2}^2
=\tfrac{9\eta_G^2}\pi\|\varphi\|_{L^2}^2$. Inserting
$\|\varphi\|_{L^2}^2=\pi(2a_0^2+a_1^2+b_1^2+\rho^2)$ and solving,
$2a_0^2+a_1^2+b_1^2\le\tfrac{9\eta_G^2}{1-9\eta_G^2}\rho^2$. With
$\eta_G\le\tfrac1{150}$ the factor is $\le\tfrac1{10}$, giving
\eqref{eq:lowmodes}.
\end{proof}

\begin{proposition}[Deficit coercivity]\label{prop:coercive}
For $\eta_G=\tfrac1{150}$, any normalised radial graph \eqref{eq:norm} with
$\|\varphi\|_{L^\infty}\le\eta_G$ and barycentre $0$ obeys
\begin{equation}\label{eq:coercive}
   \dstar(G)\ \ge\ \rho^2=\sum_{k\ge2}(a_k^2+b_k^2).
\end{equation}
\end{proposition}

\begin{proof}
Grouping \eqref{eq:quadform} by mode through \cref{lem:permode} (which already
incorporates the volume identity
$\tfrac14\int\varphi+\tfrac38\int\varphi^2=-\tfrac14\int\varphi^2$),
\begin{equation}\label{eq:Tub}
   \frac\pi8-\int_G|\sigma|^2
   =\sum_{k\ge1}D_k(a_k^2+b_k^2)-\mathcal R
   \ \ge\ \frac\pi4\,\rho^2-|\mathcal R|,
\end{equation}
discarding the (vanishing) $k=1$ term and using $D_k\ge\tfrac\pi4$ for $k\ge2$
from \eqref{eq:Dk-floor}. By \eqref{eq:Rbound}, \eqref{eq:cL2} and
\cref{lem:normmodes},
\[
   |\mathcal R|\le 10\,\eta_G\,\|\varphi\|_{L^2}^2
   \le 10\,\eta_G\cdot\tfrac{11}{10}\pi\,\rho^2=11\pi\,\eta_G\,\rho^2
   =\frac{11\pi}{150}\rho^2\le\frac\pi8\rho^2,
\]
since $\tfrac{11}{150}<\tfrac18$. Therefore
$\tfrac\pi8-\int_G|\sigma|^2\ge(\tfrac\pi4-\tfrac\pi8)\rho^2=\tfrac\pi8\rho^2$.
By \cref{lem:dual}, $T(G)\le\int_G|\sigma|^2$, so
$\tfrac\pi8-T(G)\ge\tfrac\pi8\rho^2$; dividing by $T(B_G)=\pi/8$ gives
$\dstar(G)\ge\rho^2\ge\tfrac12\rho^2$.
\end{proof}

\begin{remark}[Quantitative margin]\label{rem:numerics}
The threshold $\eta_G=\tfrac1{150}$ is dictated only by the generous
remainder constant $10$ in \eqref{eq:Rbound}; the true admissible amplitude is
far larger. A direct evaluation of the explicit admissible energy
$\int_G|\sigma|^2$ — itself a rigorous upper bound for $T(G)$ by
\cref{lem:dual}, needing no further estimate — gives the stronger
$\dstar(G)\ge\tfrac{8}\pi\cdot 0.73\,\rho^2\ge\tfrac2\pi\rho^2$ already for
$\|\varphi\|_{L^\infty}\le\tfrac1{13}$, with the infimum of
$\dstar(G)/\rho^2$ attained at the single mode $k=2$ (the ellipse), the
classical Saint--Venant extremal. We adopt the
clean lower bound \eqref{eq:coercive}, $\dstar(G)\ge\rho^2$, which already has
the sharp small-amplitude constant $1$ at $k=2$ doubled by the
volume-normalisation gain; the value of $C_G$ below uses only this bound.
\end{remark}

\subsection{The asymmetry upper bound and conclusion}

\begin{lemma}[Asymmetry bound]\label{lem:fuglede}
For a normalised radial graph \eqref{eq:norm} with
$\|\varphi\|_{L^\infty}\le\tfrac1{150}$ and barycentre $0$,
\begin{equation}\label{eq:asym}
   \AF(G)^2\ \le\ \frac{16}{\pi}\,\rho^2 .
\end{equation}
\end{lemma}

\begin{proof}
As $|G|=\pi$, the concentric unit ball is admissible in
\eqref{eq:dstarAdef}, so $\AF(G)\le\pi^{-1}|G\triangle B_1|$ and, in polar
form,
\[
   |G\triangle B_1|=\int_0^{2\pi}\Bigl|\varphi+\tfrac12\varphi^2\Bigr|\d\theta
   \le\bigl(1+\tfrac12\eta_G\bigr)\!\int_0^{2\pi}\!|\varphi|
   \le\tfrac{301}{300}\sqrt{2\pi}\,\|\varphi\|_{L^2}.
\]
Hence $\AF(G)^2\le\pi^{-2}(\tfrac{301}{300})^2 2\pi\|\varphi\|_{L^2}^2\le
\tfrac3\pi\|\varphi\|_{L^2}^2$, and by \cref{lem:normmodes}
$\|\varphi\|_{L^2}^2\le\tfrac{11}{10}\pi\rho^2$, so
$\AF(G)^2\le\tfrac{33}{10}\rho^2\le\tfrac{16}\pi\rho^2$ (as
$\tfrac{33}{10}<\tfrac{16}\pi\approx5.09$).
\end{proof}

\begin{proof}[Proof of \cref{thm:graphG}]
After the dilation and translation reducing to \eqref{eq:norm} with
barycentre $0$, \cref{prop:coercive} gives $\dstar(G)\ge\rho^2$ and
\cref{lem:fuglede} gives $\AF(G)^2\le\tfrac{16}\pi\rho^2$. Therefore
\[
   \AF(G)^2\le\frac{16}\pi\rho^2\le\frac{16}\pi\,\dstar(G)\le 48\,\dstar(G),
\]
since $\tfrac{16}\pi\approx5.09<48$. This is \eqref{eq:graphG} with
$\eta_G=\tfrac1{150}$, $C_G=48$ (any constant $\ge\tfrac{16}\pi$ works; we
report the safe value $48$).
\end{proof}

\begin{remark}[Sharpness and what is used]\label{rem:sharp}
The exponent $2$ in \eqref{eq:graphG} is sharp: for the pure mode
$\varphi=a\cos2\theta$ (a small ellipse), after normalisation
$\AF(G)^2\sim c_1 a^2$ and $\dstar(G)\sim c_2 a^2$ as $a\to0$ with
$c_1,c_2>0$, so $\AF(G)^2/\dstar(G)$ has a finite positive limit and no
exponent larger than $2$ can hold. The proof uses only the two variational
characterisations of torsion (\cref{lem:dual}; elementary), harmonicity of
$r^k e^{\mathrm i k\theta}$, the divergence theorem on the disc, Parseval, and
the volume/barycentre normalisations — \emph{no} quantitative isoperimetric
inequality, \emph{no} selection or symmetrisation, and \emph{no} boundary
regularity beyond the Lipschitz graph structure.
\end{remark}
